\section{Limitations}
\label{sec:limitations}
The attention recognition results depend on $4\times4$ pooling whose full-token control
is inconclusive. This affects both cross-tower comparisons and layer-curve shape;
sharing a pooling operation does not guarantee it distorts all layers equally.
Mean-readout agreement was established on two towers, not the entire roster.
Geometry always uses full patch tokens.

DIODE and KITTI have only 115 and 150 test images, and DIODE readouts train on just
541 images. Equal training-set size does not imply equal underfitting across models.
Bootstrap intervals resample images or pairs, not independent scenes, objects, or new
training runs; correlations among units may make them overconfident. The selected
headline comparisons are not a family-wide significance test. Layer peaks selected
from test scores are descriptive and require separate validation before operational use.

PCA controls head width but can affect representations differently, and the full matched
geometry matrix has not been rerun with deterministic PCA. Qwen3.8's $28\times28$
native grid differs from the other towers' $32\times32$ grids. A common-grid sensitivity
experiment remains missing; neither PCA nor resizing correspondence maps establishes
identical original spatial information. Kimi K3 also lacks parent forward-output and
projector-weight parity checks, weakening confidence in its extracted pipeline.

Only rectangular occlusion was measured. The results do not establish robustness to blur,
scale, lighting, or other distribution shifts. Finally, frozen probing evaluates access
under our readouts. No language model is evaluated or adapted, and we do not test
whether the observed stage orderings survive multimodal training.
